\chapter{Human review, policy, and safe export}

\section{The reviewer is part of the safety system}

Human review is not a temporary weakness to hide. It is the control that handles model uncertainty, unusual objects, unsupported classes, and consent information that pixels cannot contain.

The interface should show the normalized image, candidate overlays, class/provider labels, confidence, uncertainty, contributing evidence, and provider availability. It should allow the reviewer to:

\begin{itemize}[leftmargin=6mm]
  \item add a rectangle, polygon, or brush mask;
  \item remove an incorrect candidate;
  \item resize/expand accepted geometry;
  \item change the review label without changing raw evidence;
  \item mark an unsupported concern for manual redaction;
  \item inspect conflicts and unavailable providers; and
  \item explicitly approve the final union mask.
\end{itemize}

Every edit records user/session ID, timestamp, action, affected geometry, and reason. Undo/redo should preserve an append-only audit event stream.

\section{Consent is external state}

A face or body detector can locate a person. It cannot know whether that person agreed to a particular publication. Consent records must therefore come from user input or an external record and include scope, purpose, asset, version, timestamp, and revocation state.

The V1 policy supports conservative states such as:

\begin{itemize}[leftmargin=6mm]
  \item \file{CONSENT\_CONFIRMED};
  \item \file{CONSENT\_DENIED};
  \item \file{CONSENT\_UNKNOWN};
  \item \file{NOT\_APPLICABLE}; and
  \item \file{CONSENT\_REVOKED}.
\end{itemize}

Unknown is not equivalent to approved. If the policy requires consent for an unredacted region, unknown routes to \file{HOLD\_FOR\_CONSENT}.

\section{Policy inputs and outputs}

Policy inputs include review completion, required-provider availability, unresolved candidates/conflicts, consent state, requested output type, renderer report, and assurance report. Outputs use stable reason codes.

\begin{table}[htbp]
\centering
\caption{Representative fail-closed policy cases.}
\begin{tabularx}{\textwidth}{@{}Y L{48mm} Y@{}}
\toprule
\textbf{Condition} & \textbf{Decision} & \textbf{Reason} \\
\midrule
Review incomplete & \file{HOLD\_FOR\_REVIEW} & A proposal or manual concern is unresolved. \\
Consent required but unknown & \file{HOLD\_FOR\_CONSENT} & Pixels cannot supply consent. \\
Required face provider unavailable & \file{HOLD\_FOR\_REVIEW} & Safety evidence is missing, not negative. \\
Renderer did not cover approved mask & \file{REJECT\_EXPORT} & Output pixels do not match reviewer approval. \\
OCR or barcode attacker finds content & \file{HOLD\_FOR\_REVIEW} & Export assurance found possible leakage. \\
All required conditions pass & \file{ALLOW\_REDACTED} & Release redacted output and audit report. \\
\bottomrule
\end{tabularx}
\end{table}

\section{Destructive rendering}

V1 uses solid pixel replacement for approved regions. Blur and pixelation may be offered only as non-assured previews, not as the default privacy mechanism. The renderer must:

\begin{enumerate}[leftmargin=7mm]
  \item union all approved automatic and manual masks;
  \item verify dimensions and non-empty geometry;
  \item refuse to overwrite the source path;
  \item paint every approved pixel with the configured opaque color/pattern;
  \item encode a new JPEG or PNG from the normalized pixel array;
  \item omit source metadata and embedded previews;
  \item close and reopen the new file; and
  \item compare decoded pixels against the approved mask and rendering contract.
\end{enumerate}

The renderer report records source hash, normalized pixel hash, approved mask hash, output hash, format, dimensions, pixel count changed, encoder settings, and verification result.

\begin{warningbox}[Never trust appearance alone]
An image that looks redacted in the UI may still contain source metadata, an old thumbnail, readable blurred text, a decodable barcode, or a mask-placement error. Assurance evaluates the newly encoded file, not the preview canvas.
\end{warningbox}

\section{Independent assurance}

Assurance attacks the exported asset with tools/configurations independent from the main decision where possible:

\begin{longtable}{@{}L{33mm}L{55mm}L{57mm}@{}}
\toprule
\textbf{Check} & \textbf{Method} & \textbf{Failure behavior} \\
\midrule
\endfirsthead
\toprule
\textbf{Check} & \textbf{Method} & \textbf{Failure behavior} \\
\midrule
\endhead
Fresh decode & Reopen output with strict decoder; verify dimensions/mode & Reject corrupted or ambiguous output \\
Mask coverage & Compare approved mask with known replacement pixels & Reject any approved pixel not rendered \\
Metadata & Inspect using ExifTool independent of the writer & Reject forbidden tags/thumbnails/GPS \\
OCR & Run text detector/recognizer on output and redaction boundaries & Hold on readable sensitive text or uncertain boundary \\
Barcode & Rescan supported formats under grayscale/resize variants & Hold on any decoded payload in protected region \\
Face & Run independent/high-recall face detector & Hold when unapproved face evidence remains \\
Plate & Run plate detector/OCR attacker & Hold when unapproved plate evidence remains \\
File identity & Compare source/output hashes and paths & Reject accidental overwrite or unchanged output \\
\bottomrule
\end{longtable}

An assurance check has three states: pass, fail, uncertain/unavailable. The last two both block release. A boolean API that maps ``tool unavailable'' to pass is unsafe.

\section{Audit report}

The per-export audit report contains:

\begin{itemize}[leftmargin=6mm]
  \item source and normalized hashes without embedding private source pixels;
  \item model-bundle and threshold-profile identifiers;
  \item provider success/error/unavailable state;
  \item accepted/rejected evidence and fusion provenance;
  \item reviewer edits and approval event;
  \item consent/policy state and reason codes;
  \item renderer output hash and mask hash;
  \item every assurance result and tool version; and
  \item final decision and timestamp.
\end{itemize}

Audit reports should avoid storing recognized private strings unless required. Prefer salted hashes, redacted previews, class labels, and geometry.

\section{UI acceptance scenarios}

The reviewer workflow is complete only when the following work end to end:

\begin{enumerate}[leftmargin=7mm]
  \item no-sensitive image with zero candidates;
  \item single face requiring expansion;
  \item crowded image with missed-face manual addition;
  \item distant plate with specialist/global overlap;
  \item handwritten page with printed/handwriting uncertainty;
  \item mixed document containing signature, printed text, barcode, and metadata;
  \item provider unavailable state;
  \item assurance failure returning to the same editable review state; and
  \item approved re-export producing a new audit/output hash.
\end{enumerate}
